\section{The mean-field endpoint \texorpdfstring{$\gamma=1$}{gamma = 1}}\label{sec:meanfield}
At this endpoint the internal matrices cannot be discarded by the preceding
spectral argument. We instead view one neuron's values over all Boolean
inputs as a random vector, called its field, and track covariances of these
vectors through the layers. Conditional on the preceding layer, a
non-Gaussian row and its Gaussian comparison have the same covariance
but different distributions. The following estimate compares their log
partition functions uniformly in the field's dimension, which will be
$2^d$.

For matrices, $\|\Gamma\|_{1,\mathrm{ent}}$ denotes the sum of absolute
entries, and $\|Q\|_{\max}$ denotes the largest absolute entry. The
first norm controls tilt parameters and the second measures covariance
error; their pairing avoids a factor equal to the number of input pairs.

\begin{lemma}[A relative Lindeberg estimate]\label{lem:lindeberg}
Let $a_1,\ldots,a_N\in\R^s$ obey $\|a_j\|_\infty\leq N^{-1/2}$.
Let $\xi_j$ be independent centered variance-one variables with law $\nu$,
and let $g_j$ be independent standard Gaussians. Define
\[
 Z_\nu=\sum_j a_j\xi_j,\qquad Z_G=\sum_j a_jg_j,\qquad
 F_\Gamma(z)=\tanh(z)^T\Gamma\tanh(z),
\]
where $\|\Gamma\|_{1,\mathrm{ent}}\leq G$. Write
\[
 P(G)=24G+36G^2+8G^3,\qquad
 J_\nu=\E[|\xi|^3e^{|\xi|}]+\E[|g|^3e^{|g|}],\qquad
 c_\nu=eJ_\nu/6.
\]
If $N\geq4G^2$, then
\begin{equation}\label{eq:lindeberg}
 \left|\log\E e^{F_\Gamma(Z_\nu)}-
             \log\E e^{F_\Gamma(Z_G)}\right|
 \leq c_\nu P(G)/\sqrt N.
\end{equation}
The constant is finite for the curvature family. The conclusion requires
only the displayed moment, centering, and variance assumptions.
\end{lemma}
\begin{proof}
Using $|\tanh|,|\tanh'|\leq1$, $|\tanh''|\leq2$, and
$|\tanh'''|\leq6$, differentiation in a direction $a$ with
$\|a\|_\infty\leq N^{-1/2}$ gives
\[
 |D_aF_\Gamma|\leq2G/\sqrt N,\quad
 |D_a^2F_\Gamma|\leq6G/N,\quad
 |D_a^3F_\Gamma|\leq24G/N^{3/2}.
\]
Thus $|D_a^3e^{F_\Gamma}|\leq P(G)N^{-3/2}e^{F_\Gamma}$.
Fix the residual vector $z$ in a one-coordinate replacement and put
$f_z(t)=e^{F_\Gamma(z+at)}$ and $\lambda=2G/\sqrt N\leq1$.
The first derivative bound implies
\[
 e^{-\lambda|t|}f_z(0)\leq f_z(t)\leq e^{\lambda|t|}f_z(0).
\]
Taylor expansion through order two with integral remainder yields
\[
 |f_z(t)-f_z(0)-tf_z'(0)-t^2f_z''(0)/2|
 \leq\frac{P(G)}{6N^{3/2}}|t|^3e^{\lambda|t|}f_z(0).
\]
The first three averaged terms match for $\xi$ and $g$. The difference
of the one-coordinate expectations is therefore at most
$P(G)J_\nu f_z(0)/(6N^{3/2})$. Both expectations are at least
$e^{-1}f_z(0)$, since Jensen and variance one give
$\E e^{-\lambda|\xi|}\geq e^{-\lambda\E|\xi|}\geq e^{-1}$.
The same holds for $g$.

Averaging over any independent random residual $z$ preserves both bounds.
The inequality $|\log u-\log v|\leq|u-v|/\min\{u,v\}$ shows that
the logarithms differ by at most $c_\nu P(G)N^{-3/2}$ at this
replacement. Telescoping over all $N$ coordinates proves
\eqref{eq:lindeberg}. No inverse global partition-function bound was
used; such a bound would have introduced an unnecessary exponential in
$G$.
\end{proof}

Let $s=2^d$, $r=L-1$, and index fields and covariance matrices by all
inputs. Write
\[
 Q_{\ell,N}=\frac1N\sum_{i=1}^N h_i^\ell(h_i^\ell)^T.
\]
For a law $\mu$ on $[-1,1]^s$, define
\begin{equation}\label{eq:certificate}
 \mathcal J_G(\widehat Q;\mu)=
 \max_{\|\Gamma\|_{1,\mathrm{ent}}\leq G}
 \left\{\langle\Gamma,\widehat Q\rangle
       -\log\E_{h\sim\mu}e^{h^T\Gamma h}\right\}.
\end{equation}
Appendix~\ref{sec:sparse-dual} proves, for any law $\mu$, that an
empirical second moment from $N$ independent rows satisfies
\begin{equation}\label{eq:certificate-tail}
 \Pr\{\mathcal J_G(\widehat Q;\mu)>K\}
 \leq\exp\{q\log(2s^2+1)-N(K-1)\},\qquad q=\lceil G^2/2\rceil.
\end{equation}
It also proves that $\mathcal J_G\leq K$ supplies a quadratic
exponential tilt, of parameter norm at most $G$, whose second moment is
within $K/G$ of $\widehat Q$ in entrywise maximum norm. These facts
are distribution-free.

Let $\mu_0$ be the actual first-layer activation-field law, obtained
from independent $\nu$ weights divided by $\sqrt d$. Keep this law
exactly; its dimension need not tend to infinity. For later layers write
$\mathsf M_Q$ for the activation law of a Gaussian field of covariance
$Q$, and $\mathsf F_\Gamma(Q)$ for its second moment after tilting by
$e^{h^T\Gamma h}$. Lemma~\ref{lem:tilted-covariance} proves
\begin{equation}\label{eq:covariance-stability}
 \|\mathsf F_\Gamma(Q)-\mathsf F_\Gamma(Q')\|_{\max}
 \leq8(1+G)^2\|Q-Q'\|_{\max}.
\end{equation}
This comparison concerns an auxiliary Gaussian family; the Bayesian
learner still has non-Gaussian weights at every layer.

Conditional on the preceding activations $h_1,\ldots,h_N$, an internal
row has preactivation $N^{-1/2}\sum_j\xi_jh_j$. Apply
Lemma~\ref{lem:lindeberg} with $a_j=h_j/\sqrt N$. Its Gaussian
comparison has covariance exactly $Q=N^{-1}\sum_jh_jh_j^T$.
If $\mu_{\mathrm{act}}$ is the true conditional activation law, then
\begin{equation}\label{eq:certificate-transfer}
 \mathcal J_G(\widehat Q;\mathsf M_Q)
 \leq\mathcal J_G(\widehat Q;\mu_{\mathrm{act}})
       +c_\nu P(G)/\sqrt N.
\end{equation}
This step does not assume that the actual row law is determined by $Q$.

\subsection{A likely finite class}
The covariance comparisons now produce a finite class containing every
global-margin function outside a small parameter event. We choose the
accuracy budgets backward through depth, so that errors accumulated by
the forward reconstruction stay below the output margin. These budgets
remain polynomial only when depth is fixed.
Let $K\geq4$ and $\bar K=K+1$. Choose the budgets backward:
\begin{equation}\label{eq:budgets}
 G_r=32r\bar K^2,\qquad
 G_\ell=8G_{\ell+1}(1+G_{\ell+1})^2\quad(\ell<r),
 \qquad b_\ell=8(1+G_\ell)^2.
\end{equation}
Set $q_1=\lceil G_1^2/2\rceil$, $\ell_0=\log(2s^2+1)$, and
\begin{equation}\label{eq:meanfield-threshold}
 T(K)=\left\lceil\max\left\{1,\frac{2q_1\ell_0}{K-1},
                    4G_1^2,c_\nu^2P(G_1)^2\right\}\right\rceil.
\end{equation}
These quantities have polynomial majorants in $d,K$ at each fixed depth:
$\ell_0\leq(2d+2)\log2$ and the recursion in \eqref{eq:budgets} has
only the fixed number $r$ of steps.

\begin{lemma}[Likely proxy covariances]\label{lem:proxy}
There is a constant $c>0$, depending only on the coordinate law, such
that for all sufficiently large fixed-law $K$ and $N\geq T(K)$,
except on a parameter event of probability at most $Le^{-cNK}$, the
following hold simultaneously: $N\|W^L\|_2^2\leq K$, and there are
symmetric matrices $\Gamma_\ell$ with
$\|\Gamma_\ell\|_{1,\mathrm{ent}}\leq G_\ell$ such that the tower
\[
 Q_1^*=\frac{\E_{h\sim\mu_0}[hh^Te^{h^T\Gamma_1h}]}
                    {\E_{h\sim\mu_0}e^{h^T\Gamma_1h}},
 \qquad Q_\ell^*=\mathsf F_{\Gamma_\ell}(Q_{\ell-1}^*)\ (\ell\geq2)
\]
satisfies $\|Q_{r,N}-Q_r^*\|_{\max}\leq1/(32\bar K)$.
\end{lemma}
\begin{proof}
Apply \eqref{eq:certificate-tail} to each actual conditional row law,
including $\mu_0$ in the first layer. The threshold gives each failure
probability at most
$e^{-N(K-1)/2}\leq e^{-NK/4}$. By
\eqref{eq:certificate-transfer} and \eqref{eq:meanfield-threshold},
every internal Gaussian certificate is at most $K+1=\bar K$.
The first-layer certificate is at most $K$ and needs no comparison.
The bounded-dual lift in Appendix~\ref{sec:sparse-dual} therefore
supplies the tilt matrices with local covariance error at most
$\bar K/G_\ell$.

Let $e_\ell=\|Q_{\ell,N}-Q_\ell^*\|_{\max}$.
Then $e_1\leq\bar K/G_1$, and
$e_\ell\leq\bar K/G_\ell+b_\ell e_{\ell-1}$ for $\ell\geq2$,
by \eqref{eq:covariance-stability}. The backward budgets give
\[
 e_r\leq\sum_{\ell=1}^r\frac{\bar K}{G_\ell}
                  \prod_{j=\ell+1}^r b_j
       =\frac1{32\bar K}.
\]
For the readout, \eqref{eq:square-moment} gives
\[
 \Pr\{N\|W^L\|_2^2>K\}
 \leq\exp\{-mNK/4+(N/2)\log2\}\leq e^{-mNK/8}
\]
when $K\geq4\log2/m$. Thus $c=\min\{1/4,m/8\}$ works after
requiring $K\geq\max\{4,4\log2/m\}$.
\end{proof}

\begin{lemma}[A bounded dual readout]\label{lem:readout-dual}
Suppose $H$ is an $s\times N$ feature matrix and a readout $v$ obeys
$g_i(Hv)_i>1/2$ and $N\|v\|_2^2\leq K$. Put $Q=HH^T/N$.
There is $\beta\in\R^s$ with $\|\beta\|_1\leq2K$ and
$g_i(Q\beta)_i\geq1/2$ for every $i$.
\end{lemma}
\begin{proof}
Minimize $N\|u\|_2^2/2$ subject to $g_i(Hu)_i\geq1/2$.
The strictly feasible $v$ gives the convex optimality conditions with
multipliers $\lambda_i\geq0$:
$Nu=H^T\beta$, $\beta_i=g_i\lambda_i$, and
$\lambda_i[g_i(Hu)_i-1/2]=0$. Hence $Hu=Q\beta$ and
\[
 \tfrac12\|\beta\|_1=\beta^THu=N\|u\|_2^2\leq K.
\]
No inversion or nonsingularity of $Q$ is needed.
\end{proof}

\begin{proposition}[Mean-field aggregate mass bound]\label{prop:mf-aggregate}
For each $K$ as in Lemma~\ref{lem:proxy}, there is a finite class
$\mathcal H_K$, independent of learner width, with
$\log|\mathcal H_K|\leq\widetilde H_{L,\nu}(d,K)$ for a
polynomially bounded function at fixed depth, such that
\[
 \mu_{N,1}\!\left(\bigcup_{g\notin\mathcal H_K}U_{N,g}\right)
 \leq Le^{-cNK}\qquad(N\geq T(K)).
\]
\end{proposition}
\begin{proof}
Fix a full-margin draw in the likely set of Lemma~\ref{lem:proxy}.
Lemma~\ref{lem:readout-dual} gives $\|\beta\|_1\leq2K$ and margin
at least $1/2$ for $Q_{r,N}\beta$. Replacing $Q_{r,N}$ by $Q_r^*$
changes the output by at most $2K/(32\bar K)\leq1/16$.
Consequently $2Q_r^*\beta$ has signed margin at least $7/8$.

Proposition~\ref{prop:tilt-compress} realizes this tower within $1/8$
by a finite depth-$L$ network with width and weights polynomially bounded
in $d,\bar K,G_1$. Its first hidden layer is sampled from the true
first-layer law with the tilt; subsequent layers use Gaussian rows to
realize the auxiliary Gaussian tower. This constructs a deterministic
short description of the sign function, without changing the learner.
The realized margin is at least $3/4$. Define $\mathcal H_K$ to be all
sign functions admitting these common width and weight bounds with margin
at least $3/4$, padding by zero neurons to the common width when needed.
Lemma~\ref{lem:tilt-count} bounds this class's logarithmic cardinality
polynomially in $d,K$ at fixed depth. In particular every member of
$\mathcal H_K$ has the stated bounded implementation; the grid is used
only to count such functions. Every full-margin draw outside this class
must fail the likely event, giving the stated aggregate bound.
\end{proof}

\begin{proof}[Proof of Theorem~\ref{thm:finite-main} for $\gamma=1$]
Let $A=A_L(d,C)$ and choose
\[
 K\geq\max\{4,4\log2/m,\ c^{-1}[A+\log(6L/\varepsilon)]\}.
\]
For $N\geq T(K)$, Proposition~\ref{prop:mf-aggregate} and the target
mass give
\[
 \frac{\pi(\mathcal H_K^c)}{\pi(f)}
 \leq Le^{-N(cK-A)}\leq\varepsilon/6.
\]
The target therefore belongs to $\mathcal H_K$.
Lemma~\ref{lem:finite-posterior} gives risk at most $\varepsilon/3$
provided $n\geq(6/\varepsilon)[\widetilde H_{L,\nu}(d,K)+
\log(1/\delta)]$.
For every remaining width $C\leq N<T(K)$, the target lower bound gives
$-\log\pi(f)\leq AT(K)$. Lemma~\ref{lem:target-posterior} supplies
the same risk at sample size
$(6/\varepsilon)[AT(K)+\log(6/(\varepsilon\delta))]$.
Both $\widetilde H_{L,\nu}(d,K)$ and $AT(K)$ have polynomial
majorants in $d,C,\log(1/\varepsilon)$ at fixed depth. Taking a common
majorant and applying \eqref{eq:confidence} completes the proof.
\end{proof}
